\documentclass[../main.tex]{subfiles}

\begin{document}

\chapter*{Conclusion générale}
\addcontentsline{toc}{chapter}{Conclusion générale}

Ce projet de fin d'études illustre comment une démarche d'ingénieur rigoureuse ---
faite d'observation, de détection de problèmes et de proposition de solutions ---
peut transformer une mission initiale de modélisation en un projet à double impact.

Débuté comme une mission de modélisation stratégique pure, le projet a évolué vers
un périmètre de modélisation et de digitalisation après que nous avons identifié,
lors de l'immersion métier, l'absence d'outil de gouvernance de la donnée interne
au sein du Pôle Business Développement d'Atlantic Re. Cette détection proactive et
la proposition qui en a découlé constituent, au-delà des livrables techniques, l'apport
que nous considérons comme le plus significatif de ce stage : la capacité à dépasser
la commande initiale pour apporter une valeur ajoutée réelle, identifiée par
l'observation directe du terrain.

Sur le plan des réalisations, deux livrables opérationnels ont été produits.
L'écosystème digital de gouvernance des données interne a éliminé les sources
d'incohérence liées aux extractions manuelles et offre aux équipes un portail
structuré de consultation et d'export des KPI. Le pipeline de modélisation a produit
un panel à zéro valeur manquante sur l'ensemble des marchés africains retenus,
validé par une double procédure statistique (LOCO) et métier, sur lequel repose le
modèle SCAR en cours de finalisation.

Les perspectives naturelles de ce travail incluent l'automatisation du pipeline de
mise à jour annuelle du panel, l'intégration de nouvelles sources de données
(exposition aux risques climatiques, données télématiques pour l'automobile) et
l'extension du modèle SCAR aux branches complémentaires une fois le périmètre initial
stabilisé.

\clearpage

% ============================================================
%  RÉFÉRENCES BIBLIOGRAPHIQUES
% ============================================================

\end{document}
